% !TeX root = ../main.tex
% RESUMEN DEL PROYECTO. Normativa habitual: 2-4 páginas con los cuatro
% apartados fijos. Puede incluir alguna figura (se numera aparte):
%   \begin{center}
%     \includegraphics[width=0.7\linewidth]{generadas/ejemplo_curva}
%     \captionof{figure}{Descripción breve.}
%   \end{center}
\begin{resumenproyecto}{RESUMEN DEL PROYECTO}

\apartadoresumen{1. Introducción}
Contexto del problema, motivación y objetivos del trabajo en pocas líneas.

\apartadoresumen{2. Metodología}
Qué se ha hecho y cómo: enfoque, herramientas y decisiones principales.

\apartadoresumen{3. Resultados}
Resultados cuantitativos más relevantes. Usa cifras concretas, por
ejemplo un error del \qty{3.9}{\percent} frente al \qty{6.9}{\percent}
de la referencia.

\apartadoresumen{4. Conclusiones}
Qué se concluye y qué queda abierto.

\vspace{1em}
\noindent\textbf{Palabras clave:} \PalabrasClave

\end{resumenproyecto}
